\chapter{モデルと学習手法}

\section{モデルの入力と出力}

実運用モデルは、3台のカメラ画像、双腕の現在状態、タスク指示を同時に受け取り、両腕の関節とグリッパの将来動作を直接生成する視覚・言語・動作モデルである。本プロジェクトでは多段階タスク向けの $\pi_{0.5}$ 系列を採用した\cite{pi05_2025}。ただし本章の数値と構成は、論文上の一般論ではなく、実際の保存モデル、設定、学習履歴を根拠とする。

\figfull[.98\textwidth]{model_dataflow.pdf}{実運用モデルのデータフロー。上部と両手首の3視野、14変数のロボット状態、タスク指示から、将来50制御時刻の双腕動作を生成する。}

\begin{table}[H]
\centering
\caption{実運用モデルの入力と出力}
\begin{tabularx}{\textwidth}{>{\bfseries}p{31mm}p{34mm}X}
\toprule
項目 & 形状・規模 & 役割\\
\midrule
上部画像 & $224\times224\times3$ & 作業台全体、候補ボトル、両腕の位置関係を捉える\\
左手首画像 & $224\times224\times3$ & 左グリッパ、ボトル本体、遮蔽部分を近距離で捉える\\
右手首画像 & $224\times224\times3$ & ボトル口、キャップ、右グリッパ、回転接触を捉える\\
ロボット状態 & 離散化した14変数 & 双腕と両グリッパの現在位置を表す\\
タスク指示 & 最大200言語単位 & 実行すべき工程を指示する\\
動作出力 & $50\times14$ & 将来50制御時刻の双腕関節・グリッパ目標\\
\bottomrule
\end{tabularx}
\end{table}

\section{モデル全体構成}

学習済みモデル（学習途中で保存したモデル状態）は、画像符号化部、言語・タスク理解部、動作生成部からなる。画像は小領域へ分割して符号化され、3視野、言語、離散化したロボット状態が統合時系列モデルへ入る。動作生成部はその統合情報から将来50制御周期分の動作列を生成する。アテンション重みの記録（以下、アテンション記録）は18層、50個の動作候補、各カメラ $16\times16$ の画像格子を含む。\ev{E09,E12}

\begin{table}[H]
\centering
\caption{実運用モデルの主要設定}
\begin{tabularx}{\textwidth}{>{\bfseries}p{38mm}p{35mm}X}
\toprule
項目 & 確認値 & 説明\\
\midrule
モデル種別 & $\pi_{0.5}$ 視覚・言語・動作モデル & 連続動作生成部を使用\\
保存パラメータ数 & 838,358,468 & 51個のパラメータ配列；モデルパラメータのみ保存\\
数値精度 & bfloat16 & 学習時の演算精度\\
画像入力 & 3視野、各 $224\times224$ & 下部カメラは本番運用モデルで不使用\\
状態／実動作 & 14／14次元 & モデル内部では動作を32次元の表現領域に格納し、実機出力には14次元のみを使用\\
一度に生成する動作列 & 50制御周期 & 約1秒分の制御計画\\
推論時の反復更新回数 & 10 & ガウスノイズで初期化した候補動作を計算機内で10回更新\\
学習方式 & 全パラメータ微調整 & 当該設定では凍結なし\\
事前学習元 & $\pi_{0.5}$ 基盤モデル & 事前学習データは自社収集データとは別\\
指数移動平均 & 学習設定0.99 & 運用保存物に独立した平均パラメータは未確認\\
\bottomrule
\end{tabularx}
\end{table}

\section{動作学習の目的関数}

実演による正解動作列を $\mathbf{a}$、ランダムノイズを
$\boldsymbol{\epsilon}\sim\mathcal{N}(0,I)$ とし、
$t=0.001+0.999\,\operatorname{Beta}(1.5,1)$ を標本化する。学習時の中間動作は
\[
\mathbf{x}_t=t\boldsymbol{\epsilon}+(1-t)\mathbf{a},
\]
正しい修正方向は
\[
\mathbf{u}_t=\boldsymbol{\epsilon}-\mathbf{a}
\]
である。観測 $\mathbf{o}$ は3視野、ロボット状態、タスク指示を表し、実際の損失は
\[
\mathcal{L}_{\mathrm{flow}}
=\mathbb{E}\!\left[
\frac{1}{D}\left\|
\mathbf{v}_\theta(\mathbf{x}_t,t,\mathbf{o})-\mathbf{u}_t
\right\|_2^2\right],
\qquad D=32
\]
である。実機は内部32次元のうち14次元を利用する。\ev{E06}

\plain{正解動作をいったんノイズで崩し、どの方向へ直せば人の教示動作へ戻るかを学ばせる。実機がランダム動作を実行するのではなく、実機へ送る前に、計算機内の候補動作をノイズから段階的に更新する。}

推論は、ガウスノイズで初期化した候補 $\mathbf{x}_1$ から始め、10回の等間隔更新で $\mathbf{x}_0$ へ進める。
\[
\mathbf{x}_{t-\Delta t}
=\mathbf{x}_{t}-\Delta t\,\mathbf{v}_\theta(\mathbf{x}_t,t,\mathbf{o}),
\qquad \Delta t=0.1.
\]
最終的に50時刻分の動作列が得られる。この符号方向は確認済みの学習処理と一致する。

\section{学習設定と計算資源}

\begin{table}[H]
\centering
\caption{本番運用モデルの学習設定}
\begin{tabularx}{\textwidth}{>{\bfseries}p{37mm}p{38mm}X}
\toprule
項目 & 確認値 & 意味\\
\midrule
バッチサイズ & 256 & 1回の更新で用いる教示サンプル数\\
並列計算 & NVIDIA H200 GPU 4基 & 分散全パラメータ学習\\
データ読み込み処理 & 64 & 多数の確定版データと映像の読み込みを高速化\\
学習時の乱数初期値 & 42 & 1種類のみで、乱数初期値による変動は未評価\\
学習率 & 最大 $2.5\times10^{-5}$ & 10,000ステップで上昇後、$2.5\times10^{-6}$ へ減衰\\
最適化 & AdamW系 & 勾配裁断上限1.0\\
保存間隔 & 1,000ステップ & 現場はステップ19,000を採用\\
履歴範囲 & 0--59,990ステップ & 記録時間は約51.4時間\\
学習／検証／試験 & 学習のみ & 当該履歴に検証・試験指標なし\\
\bottomrule
\end{tabularx}
\end{table}

元設定は40,000ステップだが、履歴は59,990ステップまで継続している。学習継続は確認できる一方、40,000ステップ以後にデータ重みが変わったかは現資料から断定できない。現場運用モデルは損失最小点でも学習終点でもなく、ステップ19,000である。\ev{E08}

\section{アテンションマップの解釈範囲}

\figfull[\textwidth]{attention_real_example.png}{実行時の3視野アテンションマップ例。色の強い領域は動作生成時に重点参照した位置を表すが、この記録には成功・失敗ラベルがない。}

\figfull[.82\textwidth]{attention_camera_share.pdf}{8,223件のアテンション記録における視点別構成比。手首視点に割り当てられた重みが上部視点より相対的に大きい傾向を示す。性能への因果的寄与は遮蔽実験等で別途検証する必要がある。}

9件の実行記録一覧に8,223サンプルがあり、形式不整合や解析失敗はなかった。実行単位の平均範囲は、上部俯瞰22.9\%--25.9\%、左手首32.9\%--36.4\%、右手首37.7\%--43.4\%である。初回記録後の出力処理時間は中央値約40--55ミリ秒、95パーセンタイル約41--57ミリ秒であった。\ev{E12}

\limitbox{注意マップは、動作生成時に重点参照した画像領域を示すだけである。キャップ有無、ボトル向き、成功・失敗ラベル、遮蔽対照がないため、空回しや逆向き把持の原因を直接証明するものではない。}

\section{運用保存モデルの監査}

運用モデルの保存状態はパラメータのみを含む分割形式で、ステップ19,000の保管場所に保存されている。51個のパラメータ配列、838,358,468パラメータ、14次元の状態・動作に対する平均、標準偏差、第1・第99パーセンタイル統計を読み取れた。学習再開用の調整情報と独立した移動平均パラメータは含まれない。\ev{E09}

\begin{table}[H]
\centering
\caption{運用保存モデルの確認結果}
\begin{tabularx}{\textwidth}{>{\bfseries}p{41mm}p{31mm}X}
\toprule
確認項目 & 結果 & 解釈\\
\midrule
保存情報の読み込み & 可能 & モデル構造と統計を読み取り専用で解析\\
保存形式 & モデルパラメータのみ & 学習再開に必要な最適化状態は含まれない\\
配列数／総パラメータ数 & 51／838,358,468 & 本番運用モデルの構造と一致\\
学習ステップ & 19,000を複数資料で確認 & 運用設定、保存先、学習系譜を照合\\
正規化統計 & あり & 状態・動作各14次元、平均・標準偏差・分位数\\
独立平均パラメータ & 未確認 & 学習設定には係数があるが保存物に別複製なし\\
現行構造との整合 & 保存情報と形状が一致 & 明確な形状不一致は未検出\\
\bottomrule
\end{tabularx}
\end{table}
